Un parallamps és una barra metàl·lica vertical que atrau i dirigeix grans descàrregues de corrent cap a terra. La gran majoria de llamps núvol-terra són negatius, és a dir, són transferències de càrrega negativa del núvol cap a terra. En el moment de la descàrrega es crea un camp magnètic al voltant del parallamps que podem equiparar al creat per un fil de corrent infinit. El corrent màxim que pot assumir un parallamps és d'uns 100~kA.

\begin{apartat}{1,25}
Calculeu el camp magnètic màxim que pot crear el parallamps a una distància de 10~cm. Feu un dibuix esquemàtic del parallamps indicant el sentit del moviment dels electrons, la intensitat de corrent i tres línies de camp magnètic. Justifiqueu el sentit de les línies de camp.
\end{apartat}

\begin{apartat}{1,25}
Representeu gràficament, en la quadrícula de sota, el mòdul d'aquest camp magnètic màxim en funció de la distància $r$ al parallamps en l'interval següent: $10\ \mathrm{cm}\le r\le 50\ \mathrm{cm}$. Suposem que hi ha un electró que en el moment de la descàrrega es troba a 10~cm del parallamps i que té una velocitat de $10^{3}\ \mathrm{m/s}$ paral·lela al parallamps i cap a terra. Calculeu el mòdul de la força que crea el camp magnètic sobre l'electró i justifiqueu-ne la direcció i el sentit.
\end{apartat}

\figura[0.5]{fig1.png}

\dades{%
  El mòdul del camp magnètic creat per un fil infinit per on circula un corrent $I$ a una distància $r$ del fil és $B = \dfrac{\mu_0\, I}{2\pi\, r}$.\\
  $\mu_0 = 4\pi\times 10^{-7}\ \mathrm{T\,m\,A^{-1}}$.\\
  $|e| = 1,602\times 10^{-19}\ \mathrm{C}$.}
